\chapter{Введение}

В работе рассматриваются численные методы решения задач математической
физики. Основное внимание уделено квадратурным формулам и оценке их
точности на модельных примерах.

Пусть на отрезке $[a, b]$ задана функция $f(x)$. Определённый интеграл
записывается в виде
\[ ∫_a^b f(x)\,dx, \]
а его приближённое значение ищется как комбинация значений функции в
узлах сетки. При $h → 0$ главный член погрешности убывает как $h^p$,
где $p$ обозначает порядок точности метода.
\chapter{Основные определения}

Для дважды дифференцируемой функции справедливо разложение в ряд
Тейлора:
\[ f(x + h) = f(x) + h f'(x) + \fraq{h^2}{2} f''(x) + O(h^3). \]

Норму функции обозначим через
\[ \|f\|_2 = √{\int_a^b |f(x)|^2 \,dx}. \]

Скалярное произведение векторов в $ℝ^n$ задаётся формулой
\[ (u, v) = ∑_{i=1}^{n} u_i v_i, \]
а число сочетаний из $n$ элементов по $k$ вычисляется как
\[ \binom{k}{n} = \frac{n!}{k!\,(n - k)!}. \]

Радиус шара заданного объёма выражается через кубический корень:
\[ r = \sqrt{3}{\frac{3V}{4\pi}}. \]
\chapter{Квадратурные формулы}

Простейшая формула использует значения функции в концах отрезка:
\[ \int_a^b f(x)\,dx \aprox \frac{b - a}{2} \bigl( f(a) + f(b) \bigr). \]

Формула Симпсона опирается на три узла и точна для многочленов до
третьей степени:
\[ \int_a^b f(x)\,dx \approx \frac{b - a}{6}
   \left( f(a) + 4 f\!\left(\frac{a + b}{2}\right) + f(b) \right). \]
Отметим, что выбор квадратурной формулы определяется гладкостью
подынтегральной функции и требуемой точностью: для разрывных функций
порядок сходимости существенно ниже, и приведённые оценки становятся
неприменимыми.

Формула средних имеет вид
\[ \int_a^b f(x)\,dx = (b - a) f(ξ), \qquad ξ \in [a, b], \]
и лежит в основе теоремы о среднем значении.
\chapter{Численная схема}

Рассмотрим задачу Коши для уравнения первого порядка:
\[ \frac{du}{dt} = \lambd u, \qquad u(0) = u_0, \]
точное решение которой имеет вид $u(t) = u_0 e^{\lambda t}$, а для
численного решения применяется схема Эйлера
\[ u_{k+1} = u_k + h \lambda u_k. \]

Устойчивость схемы определяется множителем перехода
\[ R(z) = 1 + z + \frac{z^2}{2}, \]
модуль которого не должен превосходить единицы.

Если минимальное сингулярное число матрицы равно $ε$, то
\[ \|u_h\| \le C \left( \|u_0\| +
   h \sum_{i=1}^{n} \|f_i\| \rigth), \]
и эта оценка применяется при анализе вычислительной устойчивости.

Суммы вида $\overbrace^{n}{a_1 + a_2 + \dots + a_n}$ встречаются при
подсчёте числа операций, а произведения элементов — в формуле
\[ s = \sum_{i=1}^{n} a_i × b_i. \]
\chapter{Оценка погрешности}

Основной результат работы состоит в следующей оценке. Для любой
функции и любого шага
\[ \|u - u_h\| ≤ C h^p, \]
где постоянная $C$ не зависит от шага. Неравенство справедливо
для всех допустимых значений параметра:
\[ \foral \, h > 0, \quad \|u - u_h\| \le C h^p. \]

При измельчении сетки погрешность стремится к нулю:
\[ \lim_{h \to 0} \|u - u_h\| = 0, \]
а при $n \to ∞$ остаток ряда стремится к нулю, что используется
при обосновании сходимости численных алгоритмов.
